% !TeX root = ../papers/nc-p3.tex
\section{四次元族の固有な四元群・周期と同型判定}
\label{sec:klein-generic-isomorphisms}
前節の第二解の族で、三つの点配置を$X_E,X_F,X_G$と書く。

\begin{theorem}\label{thm:klein-generic-period-four}
幾何学的一般点で各配置の射影安定群は指定された四元群に等しく、三配置は互いに
射影同値でない。この標準AS窓の最小周期は4である。
さらに明示的な四つの有理不変量を高々十六通りの文字置換で比較すれば、
この条件を持つ開集合上の正則窓の添字保持同型が判定できる。
\end{theorem}
\begin{proof}
前節の素数12619の入力で、偶二十点の両評価行列は階数3である。
核の余因子から各々の奇配置を計算する。
各配置には三点直線がなく、四点以上の平面が48枚あり、それぞれ丁度四点を持つ。
このうち八平面を通す点が四つ、十平面を通す点が十六ある。
前者の四点は線型に独立で、その四つの座標面から外れる点もある。
従って配置間の射影同型は、四点の置換と、固定した第五点の行き先とで一意に決まる。
これらを全て作用させて二十点の一致を調べると、$(E,F,G)$の順の
射影同型の個数行列は$\operatorname{diag}(4,4,4)$である。
全点が素体上で第五点までが射影枠なので、これは代数閉包上でも同じ個数である。
\path{computations/verify_klein_generic_isomorphisms.py}がこの有限計算を行う。

この比較は一般点での上界を与える。有限エタールな順序付け被覆で二十点に番号を付け、
各置換$\pi$について$D x_i\wedge y_{\pi(i)}=0$を
$[D]\in\mathbb P(\operatorname{Mat}_4)$に課す。解の像は射影性から閉である。
特殊点で階数$1,2,3$の行列の核が含む点はそれぞれ高々$4,2,1$個で、
像空間が含む点は高々$1,2,4$個である。この和は20未満なので特異解はない。
番号は両奇配置にも核で運べる。実際、特殊点で全評価核は直線で相異なる奇点を与え、
偶側の有限エタール被覆上の階数3性は近傍で保たれる。両行列から得る偶点スキームの
一致は一般正則点で前述の通りであり、ある開集合へ延びる。より直接には
第一窓の各零点対で$\gamma(x,-)$は$\beta(-,f)^\perp$に入る。
特殊点で$\beta(-,f)$も階数3であることは、薄い次数が全成分非零で交差行列が
可逆なこの例では直接評価でも確認できる。従って近傍で核の点列が定まる。
有限個の置換を比較することで、一般点では射影同型の個数は上で数えた値以下となる。
指定された対角四元群は各配置を保つため各対角成分は丁度4である。
周期2なら一方の奇配置を他方に移す線型同型が必要で、周期1でも同様である。
四窓表示による周期4から最小性を得る。

従って窓の同型は偶配置の唯一の四元群を保つ。奇配置の変換も、偶配置の作用と
そのグラフ上で整合するため群上の同じ自己同型を誘導する。各配置の点ごとの
射影固定群は自明である。四文字の置換はある$\phi\in\operatorname{Aut}(K)$で
$ i\mapsto\phi(i)+a,\ j\mapsto\phi(j)+b,\ l\mapsto\phi(l)+b$
となる。実際、射影表現の持ち上げの差は共通の一文字の平行移動である。
偶二位置の変換は反対階数10から一致し、テンソルが全次数零なので奇二位置の移動量は
等しい。厚い次数の組を保つ条件は$\phi(1)=1$, $a+b\in\{0,1\}$である。
よって候補は$2\cdot4\cdot2=16$通りとなる。

文字を固定した後は、$U,V$が全成分非零で$S_{00}\ne0$の範囲を取る。
薄い成分の変換は$e_i f_j t^{\pm1}b_{ij}$であり、次数2に正、次数3に負を取る。
$e_i,f_j$は偶・奇座標の倍率で、$t$は二つの薄い次数の逆倍率である。
$\rho_U=U_{00}U_{11}/(U_{01}U_{10})$と置くと$t^4=1/\rho_U$で
$U$の行列積比が1になる。二つの行列の行・列の倍率を独立に選び、
$U$の成分を全て1、$V$の右下以外を1にすると、その右下は
$z=\rho_V/\rho_U$となる。この倍率が$S$の第1行と第0行に与える比と
第1列と第0列に与える比は、それぞれ
\[
 L=\frac{U_{01}V_{00}}{U_{11}V_{01}},\qquad
 R=\frac{U_{00}V_{01}}{U_{10}V_{00}}.
\]
実際、行比は$e_1 f_1/(e_0 f_0)$、列比は$e_1 f_0/(e_0 f_1)$である。
二組の間の残る共通倍率で$S_{00}=1$にできる。従って正規形の四係数は
\[
 \left(z,\ R\frac{S_{01}}{S_{00}},\ L\frac{S_{10}}{S_{00}},
               \ LR\frac{S_{11}}{S_{00}}\right).
\]
どの四乗根を選んでもこの表示は変わらない。逆に四係数は規格化した$U,V,S$を与え、
第二解は有理式で決まる。以上の分母と第二解の構成が十六置換の全てで有効な
開集合を取れば、置換して四係数を比較する必要十分条件となる（各置換は族の
一般点を保つ。厚い行列の各組内ブロックは一般に全成分非零で可逆である）。
正則性の主張は前節の幾何学的一般点、同型の判定は条件を満たす全正則入力に関する。
\end{proof}
